% БҮЛЭГ 5 — ЗАГВАР ХӨГЖҮҮЛЭЛТ БА PIPELINE

\chapter{Загвар хөгжүүлэлт ба дамжлага}
\label{Chapter5}

Энэхүү бүлэгт MN-BERT суурьтай хоёр шатлалт системийн сургалтын дамжлагыг
нарийвчлан баримтжуулна. Үүнд загварын бүтэц, гиперпараметрийн оновчлол,
ангиллын тэнцвэржүүлэлт, туршилтын давтагдах байдлын нөхцөлийг багтаасан.
Ийнхүүлснээр дараагийн бүлэгт (бүлэг~\ref{Chapter6}) танилцуулагдах
туршилтын үр дүнг үнэмшилтэйгээр ойлгуулах суурийг тавина.

%-------------------------------------------------------------------------------
\section{MN-BERT загварын архитектурын дэлгэрэнгүй тайлбар}

MN-BERT нь Монгол хэлний corpus дээр pre-train хийгдсэн BERT--base архитектур
бөгөөд \texttt{bert-base-mongolian-cased} хувилбарыг суурь загвар болгон
ашигласан. Уг загвар нь 12 Transformer энкодер давхарга (encoder layer), 12
attention head, далд давхаргын хэмжээ 768, нийт ойролцоогоор 110 сая параметртэй.
Оролтын хамгийн их урт нь 512 token бөгөөд ангиллын даалгаварт $256$ token-ы
урт ашиглахаар тохируулсан. Сургалтын олонлог дээрх токений урт нь медиан $29$,
95-р хувьцаа $96$, 99-р хувьцаа $135$, дээд тал нь $374$ token байсан тул $256$
хязгаар нь урт сэтгэгдлийг (ялангуяа Бүтээлч шүүмжлэл) бараг таслахгүйгээр
багтаах боломжийг олгоно (бүлэг~\ref{Chapter4}-ийн уртын шинжилгээтэй нийцнэ).

Ангиллын даалгаварт зориулан MN-BERT-ийн дээр дараах давхаргуудыг нэмсэн:

\begin{enumerate}
  \item \textbf{Tokenization давхарга:} SentencePiece-т суурилсан subword tokenizer
    нь Монгол хэлний agglutinative морфологийг алдаагүй байдлаар
    кодолно. Оролтын текст эхэнд \texttt{[CLS]}, төгсгөлд \texttt{[SEP]} тэмдэгт
    нэмэгдэнэ.
  \item \textbf{Transformer encoder stack:} 12 давхаргын гаралт нь контекст мэдрэмжтэй
    vector илэрхийлэл үүсгэнэ. Энэ хэсэг нь pre-train хийгдсэн жингүүдийг үндэс
    болгон нарийн тохируулга хийгдэнэ.
  \item \textbf{[CLS] pooling давхарга:} Төгсгөлийн encoder давхаргын
    \texttt{[CLS]} байрлалын далд вектор нь өгүүлбэрийн нэгдсэн дүрслэл болж
    ангиллын толгойд (classification head) дамжина.
  \item \textbf{Ангиллын толгой (classification head):} Нэг бүрэн холбосон
    (fully connected) давхарга нь pooled vector-ыг ангиллын логитууд (logits)
    руу шугаман хувиргана. Stage~1-ийн хувьд гаралтын хэмжээ 3 (Эерэг,
    Саармаг, Сөрөг), Stage~2-ийн хувьд 2 (Хортой сөрөг, Бүтээлч шүүмжлэл).
  \item \textbf{Softmax давхарга:} Логитуудыг магадлалын тархалтад хувиргана.
    Инференс (inference) үедээ \texttt{argmax} ашиглан эцсийн ангиллыг тодорхойлно.
\end{enumerate}

Stage~1 болон Stage~2 нь тус тусдаа нарийн тохируулга хийгдсэн бие даасан загварууд
бөгөөд мэдээлэл солилцохгүй, каскад (cascade) хэлбэрээр холбогдсон. Энэ нь
ангилал тус бүрт өөрийн тодорхой даалгаварт тохируулан оновчлох боломж олгодог.

MN-BERT-ийн нэг Transformer энкодер блокийн дотоод бүтэц болон нарийн тохируулгатай
ангиллын толгойн дэлгэрэнгүй диаграмыг зураг~\ref{fig:bert-detail}-д харуулав.

\begin{figure}[H]
\centering
\begin{tikzpicture}[
  mainbox/.style={draw, rounded corners=4pt, minimum width=7.0cm,
                  minimum height=1.20cm, text width=6.6cm,
                  align=center, font=\small, line width=0.8pt},
  arr/.style={->, >=Stealth, thick},
  resarr/.style={->, >=Stealth, dashed, gray!60, line width=0.9pt}
]

% Node spacing: 2.00 cm centre-to-centre → 0.80 cm clear gap between boxes
% Each mainbox half-height = 0.60 cm  (minimum height 1.20 cm)

% ── INPUT EMBEDDING ──────────────────────────────────────────────────────
\node[mainbox, fill=gray!12] (inp) at (0, 0)
  {Input Embeddings \quad ($\mathbb{R}^{768}$)\\[2pt]
   {\scriptsize Token + Position + Segment}};

% ── TRANSFORMER BLOCK BORDER ─────────────────────────────────────────────
% mha.north = -2.00+0.60 = -1.40  →  block top = -0.80
% an2.south = -8.00-0.60 = -8.60  →  block bottom = -8.90
\draw[dashed, rounded corners=6pt, fill=blue!3, draw=blue!40, line width=1pt]
  (-3.8,-0.80) rectangle (3.8,-8.90);
\node[font=\footnotesize\itshape, text=blue!60, anchor=north west]
  at (-3.75,-0.86) {Transformer encoder block ($\times 12$)};

% ── MULTI-HEAD ATTENTION ─────────────────────────────────────────────────
\node[mainbox, fill=blue!14] (mha) at (0,-2.00)
  {Multi-Head Self-Attention\\[2pt]
   {\scriptsize 12 heads, $d_k = 64$}};

% ── ADD & LAYER NORM 1 ───────────────────────────────────────────────────
\node[mainbox, fill=yellow!22] (an1) at (0,-4.00)
  {Add \& Layer Norm\\[2pt]
   {\scriptsize Residual: Input + Attention output}};

% ── FEED-FORWARD NETWORK ─────────────────────────────────────────────────
\node[mainbox, fill=blue!8] (ffn) at (0,-6.00)
  {Feed-Forward Network\\[2pt]
   {\scriptsize $768 \to 3072 \to 768$, GeLU}};

% ── ADD & LAYER NORM 2 ───────────────────────────────────────────────────
\node[mainbox, fill=yellow!22] (an2) at (0,-8.00)
  {Add \& Layer Norm\\[2pt]
   {\scriptsize Residual: Norm$_1$ + FFN output}};

% ── BLOCK OUTPUT ─────────────────────────────────────────────────────────
\node[mainbox, fill=blue!5] (bout) at (0,-10.00)
  {Block Output ($\mathbb{R}^{768}$)\\[2pt]
   {\scriptsize $\longrightarrow$ input of next block}};

% ── [CLS] TOKEN POOLING ──────────────────────────────────────────────────
\node[mainbox, fill=green!18] (cls) at (0,-12.00)
  {\texttt{[CLS]} Token Vector ($\mathbb{R}^{768}$)\\[2pt]
   {\scriptsize sentence-level representation}};

% ── LINEAR CLASSIFICATION HEAD ───────────────────────────────────────────
\node[mainbox, fill=orange!22] (fc) at (0,-14.00)
  {Linear: $768 \to C$ \quad (Focal Loss)\\[2pt]
   {\scriptsize classification logits}};

% ── SOFTMAX OUTPUT ───────────────────────────────────────────────────────
\node[mainbox, fill=red!14] (sm) at (0,-16.00)
  {Softmax $\to$ predicted class\\[2pt]
   {\scriptsize $C \in \{2,\,3,\,4\}$ ангилал}};

% ── MAIN ARROWS ──────────────────────────────────────────────────────────
\draw[arr] (inp)  -- (mha);
\draw[arr] (mha)  -- (an1);
\draw[arr] (an1)  -- (ffn);
\draw[arr] (ffn)  -- (an2);
\draw[arr] (an2)  -- (bout);
\draw[arr] (bout) -- (cls);
\draw[arr] (cls)  -- (fc);
\draw[arr] (fc)   -- (sm);

% ── RESIDUAL CONNECTIONS (explicit coords, right side) ───────────────────
% Residual 1: inp (y=0) → an1 (y=-4.00), offset x=+0.70 → right edge at x=4.20
\draw[resarr] (inp.east) -- ++(0.70,0) -- ++(0,-4.00) -- (an1.east);
\node[font=\scriptsize, text=gray!55, anchor=west] at (3.40,-2.00) {Residual 1};
% Residual 2: an1 (y=-4.00) → an2 (y=-8.00), offset x=+1.00 → right edge at x=4.50
\draw[resarr] (an1.east) -- ++(1.00,0) -- ++(0,-4.00) -- (an2.east);
\node[font=\scriptsize, text=gray!55, anchor=west] at (3.70,-6.00) {Residual 2};

\end{tikzpicture}
\caption{MN-BERT нарийн тохируулгатай загварын архитектур: нэг Transformer
  encoder block-ийн бүрэн бүтэц (тасархай хайрцаг) болон ангиллын толгой.
  Residual холболтыг баруун талд L хэлбэрийн сумаар үзүүллэв.}
\label{fig:bert-detail}
\end{figure}

%-------------------------------------------------------------------------------
\section{Stage 1 ба Stage 2 тус тусын сургалтын дамжлага}

Stage~1 болон Stage~2 хоёулаа PyTorch 2.x framework, HuggingFace Transformers
4.x ашиглан сургагдсан. Сургалтын үйл явцыг нэгдмэл дамжлагад зохион байгуулсан
бөгөөд түүнийг хүснэгт~\ref{tab:pipeline-flow}-д багцалсан.

\begin{table}[H]
\centering
\caption{Сургалтын дамжлагын үе шатуудын урсгал}
\label{tab:pipeline-flow}
\fittable{%
\begin{tabular}{@{}clp{8cm}@{}}
\toprule
\textbf{№} & \textbf{Үе шат} & \textbf{Тодорхойлолт} \\
\midrule
1 & Өгөгдөл ачаалалт       & CSV файлаас pandas ашиглан унших, train/val/test хуваалтаар ангилах \\
2 & Tokenization           & \texttt{AutoTokenizer.from\_pretrained} дуудаж (\texttt{bert-base-mongolian-cased}) загвараар токенчлох, \texttt{max\_length}=256 \\
3 & PyTorch Dataset        & \texttt{input\_ids}, \texttt{attention\_mask}, \texttt{labels} tensor-уудыг боловсруулах \\
4 & DataLoader             & Batch болгон ачаалах, \texttt{shuffle=True} train-д, \texttt{shuffle=False} val/test-д \\
5 & Forward pass           & \texttt{outputs = model(input\_ids, attention\_mask, labels)}, Focal Loss ($\gamma{=}2.0$, ангиллын жинтэй) тооцох \\
6 & Backward pass          & \texttt{loss.backward()} ба градиентын clipping (\texttt{max\_norm=1.0}) \\
7 & Optimizer step         & AdamW optimizer, weight decay зохицуулалттай \\
8 & LR scheduler           & Linear warmup (10\% алхамд), дараа нь шугаман буурах \\
9 & Валидаци               & Epoch тутам параметр тохируулгын өгөгдлийн багцаар F1-macro үнэлэх, \texttt{best} жинг хадгалах \\
10 & Чекпойнт              & \texttt{best\_model.pt} файлд жинг хадгалах (early stopping-тай) \\
\bottomrule
\end{tabular}}
\end{table}

Stage~1-ийн хувьд оролтын өгөгдөл нь Эерэг / Саармаг / Сөрөг гэсэн гурван
ангилалтай, гаралт нь 3 логиттой cross-entropy loss. Stage~2-ийн хувьд зөвхөн
Сөрөг гэж тэмдэглэгдсэн жишээнүүд ашиглагдана. Оролт нь хоёртын (binary)
ангилалтай (Хортой сөрөг ба Бүтээлч шүүмжлэл). Хоёр шат нь өөр өөр өгөгдлийн багц дээр
сургагдах тул тус тусын ``best'' жингээ хадгалж, инференс үед цувуулан дуудна.

%-------------------------------------------------------------------------------
\section{Оновчтой шийдэл: Нэг шатлалт (Flat) загварт шилжсэн нь}

Хоёр шатлалт каскад архитектурыг турших явцад гарсан үр дүн болон алдааны
дүн шинжилгээ (бүлэг \ref{Chapter6}-д дэлгэрэнгүй үзүүлсэн) нь нэгэн чухал
инженерийн сорилтыг илрүүлсэн. Энэ нь \textbf{Алдааны дамжуулалт (Error
Propagation)} буюу Stage~1 дээр гарсан өчүүхэн алдаа Stage~2-ын үр дүнг
бүрэн боломжгүй болгох bottleneck үүсгэсэн явдал юм.

Энэхүү олдвор дээр үндэслэн судалгааны чиглэлийг өөрчилж, дөрвөн
ангиллыг нэгэн зэрэг таамаглах \textbf{Нэг шатлалт (Flat) MN-BERT} загварыг
эцсийн шийдэл болгон сонгосон. Flat загварын давуу тал нь:

\begin{enumerate}
    \item \textbf{Алдааны нийлбэрээс сэргийлэх:} Бүх ангиллын хил хязгаарыг
    нэгэн зэрэг сурснаар Stage~1-ийн хатуу шүүлтүүрээс хамааралгүй болно.
    \item \textbf{Focal Loss ашиглалт:} Классуудын тэнцвэргүй байдлыг шийдэхэд
    Cross-Entropy-ийн оронд Focal Loss ашигласнаар ховор ангилалын (Хортой сөрөг)
    жишээнд загварын төвлөрлийг нэмэгдүүлсэн.
    \item \textbf{Инференс хурд:} Хоёр загвар цувуулан ачаалах шаардлагагүй тул
    нэг дамжуулалтаар хариу өгөх ба инференсийн хугацаа мэдэгдэхүйц богино.
\end{enumerate}

Ийнхүү бидний эцсийн систем нь Focal Loss болон Cosine Warmup тохируулга бүхий
Flat MN-BERT загвар дээр суурилсан бөгөөд энэ нь туршилтын хамгийн өндөр
гүйцэтгэлийг үзүүлсэн юм.

%-------------------------------------------------------------------------------
\section{Гиперпараметрийн тохиргоо ба оновчлол}

Гиперпараметрийн сонголтыг бүлэг~\ref{Chapter3}-т танилцуулсан тохиргооны
үндсэн дээр хийж, Stage~1 ба Stage~2-д тус тусад нь validation F1-macro
метрикээр сонгож оновчлов. Эцсийн сонголтыг хүснэгт~\ref{tab:hyperparams-mnbert}-д
жагсаав.

\begin{table}[H]
\centering
\caption{Гиперпараметрийн эцсийн тохиргоо. Эцсийн тайлагнасан систем нь
  \textbf{Нэг шатлалт (Flat)} загвар; Stage~1/Stage~2 баганууд нь харьцуулалтад
  ашигласан хоёр шатлалт хувилбарынх.}
\label{tab:hyperparams-mnbert}
\fittable{%
\begin{tabular}{@{}lccc@{}}
\toprule
\textbf{Гиперпараметр} & \textbf{Нэг шатлалт (эцсийн)} & \textbf{Stage 1} & \textbf{Stage 2} \\
\midrule
Learning rate              & $3\times10^{-5}$ & $3\times10^{-5}$ & $2\times10^{-5}$ \\
Batch size                 & 16   & 16   & 16   \\
Epochs (Best)              & 8    & 6    & 5    \\
Warmup ratio               & 0.10 & 0.1  & 0.1  \\
Weight decay               & 0.01 & 0.01 & 0.01 \\
Dropout                    & 0.15 & 0.1  & 0.1  \\
Max sequence length        & 256  & 128  & 128  \\
Loss function              & Focal Loss ($\gamma{=}2.0$) & Жинтэй CE & Жинтэй CE \\
Тэнцвэржүүлэлт             & $\alpha$ жин + WeightedRandomSampler & Жин & Жин \\
Optimizer                  & AdamW & AdamW & AdamW \\
LR scheduler               & Cosine + Warmup & Cosine + Warmup & Cosine + Warmup \\
Early stopping patience    & 3 & 3 & 3 \\
Random seed                & 42 & 42 & 42 \\
\bottomrule
\end{tabular}}
\end{table}

Learning rate-ийн хувьд $\{2\!\times\!10^{-5},\; 3\!\times\!10^{-5},\; 5\!\times\!10^{-5}\}$
гэсэн утгуудыг grid search аргаар туршин үзэж, хамгийн сайн validation F1-ыг
өгсөн утгыг сонгов. Batch size-ыг $\{16, 32\}$ хүрээнд туршихад 12GB VRAM-ийн
хязгаарлалтаас шалтгаалан batch size 16 нь сургалтын явцад хамгийн тогтвортой үр дүн өгөх нь тодорхой болсон.

%-------------------------------------------------------------------------------
\section{Сургалтын дэд бүтэц ба орчин}

Загваруудыг сургах болон турших ажлыг локал тооцооллын орчинд гүйцэтгэсэн.
Техник хангамжийн хувьд \textbf{AMD Radeon RX 9070 GRE 12GB} график картыг
ашигласан бөгөөд \texttt{torch-directml} санг ашиглан DirectX-ээр дамжуулан
GPU тооцооллыг идэвхжүүлсэн.

\begin{itemize}
    \item \textbf{GPU:} AMD Radeon RX 9070 GRE (12GB VRAM)
    \item \textbf{CPU:} Intel Core i7-12700K
    \item \textbf{RAM:} 32GB DDR5
    \item \textbf{Framework:} PyTorch 2.4.1 (with DirectML)
    \item \textbf{OS:} Windows 11
\end{itemize}

%-------------------------------------------------------------------------------
\section{Ангиллын тэнцвэржүүлэлт}

Бүлэг~\ref{Chapter4}-т үзүүлсэнчлэн, өгөгдлийн ангиллын тархалт нь тэнцвэргүй
байдлыг илэрхийлдэг: Бүтээлч шүүмжлэл (43.56\%) болон Саармаг (30.26\%) ангиллууд
давамгайлж, Эерэг (6.55\%) ангилал хамгийн цөөхөн дээжтэй. Уг тэнцвэргүй байдал нь цөөнх ангиллын F1 оноо
гүйцэтгэлийг доройтуулж болзошгүй тул загвар тус бүрд тохирсон тэнцвэржүүлэлт хийв.

\textbf{Суурь загваруудын хувьд (LR, NB, SVM):} \texttt{imbalanced-learn}
сангийн \texttt{SMOTE} (Synthetic Minority Over-sampling Technique) алгоритмыг
TF-IDF vector-ын дээр хэрэглэв. SMOTE нь цөөнх ангиллын жишээг $k$-хамгийн
ойрын хөршийнх нь интерполяциар синтетикээр үржүүлнэ. Хэрэгжүүлэлтэд
$k\!=\!5$ утгыг ашиглав.

\textbf{Эцсийн нэг шатлалт MN-BERT загварын хувьд:} BERT-д шууд SMOTE хэрэглэх нь
эмбэддингийн орон зайд семантикийг алдагдуулах эрсдэлтэй тул \textbf{Focal Loss}
($\gamma = 2.0$)-ийг ангиллын урвуу давтамжийн жин $\alpha$-тай хослуулан
ашиглав. Жин $\alpha$ нь:
\[
  \alpha_c = \frac{N}{K \cdot n_c}
\]
энд $N$ --- нийт жишээний тоо, $K$ --- ангиллын тоо, $n_c$ --- $c$-р
ангиллын жишээний тоо (тооцоолсон утга: Эерэг $3.707$, Саармаг $0.835$,
Бүтээлч $0.570$, Хортой $1.286$). Нэмж \texttt{WeightedRandomSampler}-ээр
сургалтын batch-д цөөнх ангиллыг илүү давтамжтай түүвэрлэв. Focal Loss-ийн
$(1-p_t)^{\gamma}$ хүчин зүйл нь амархан ангилагдах жишээний нөлөөг бууруулж,
хүнд (цөөнх) жишээнд төвлөрлийг нэмэгдүүлнэ.

Эдгээр аргын нөлөөг бүлэг~\ref{Chapter6}-ийн F1-macro үр дүн болон өгөгдлийн
чанаржуулалтын явц (Хүснэгт~\ref{tab:dq-progression})-аар үнэлэв.

%-------------------------------------------------------------------------------
\section{Gradio интерфейс ба дамжлагын нэгтгэл}

Эцсийн системийн хэрэглэгчийн интерфейсийг Gradio framework дээр бүтээсэн бөгөөд
гурван таб бүхий бүтэцтэй.

\begin{enumerate}
  \item \textbf{Ангилах таб (гол таб):} Хэрэглэгч нэг текст эсвэл олон текст
    (batch) оруулж, радио товч ашиглан хоёр архитектурын аль нэгийг сонгоно.
    \begin{itemize}
      \item \textit{Flat} (нэг шатлалт, нарийвчлал $\mathbf{83.26\%}$):
        \texttt{best\_mnbert\_sota\_corrected\_model} загвар текстийг нэг алхмаар
        4 ангилалд (Эерэг / Саармаг / Хортой сөрөг / Бүтээлч шүүмжлэл) шууд
        ангилна. Оролт нь сургалтын нөхцөлтэй яг нийцэхийн тулд
        \texttt{[news.mn]~normalize\_flat(text)} хэлбэрт хувиргагдаж
        \texttt{max\_length\,=\,256}-аар токенчлогдоно.
      \item \textit{Two-stage} (нарийвчлал $78.61\%$): Stage~1 загвар эхлээд
        Эерэг / Саармаг / Сөрөг гэж 3 ангилдаг; Сөрөг бол Stage~2 загвар
        Хортой сөрөг эсвэл Бүтээлч шүүмжлэлд нарийвчлан ялгана
        (\texttt{max\_length\,=\,128}).
    \end{itemize}
    Эцсийн ангилал, зэрэглэлийн итгэлцлийн хувь (softmax гаралт),
    confidence баар диаграм болон batch CSV татах боломж харагдана.
    ``Санамсаргүй жишээ'' товч нь туршилтын хэсгээс санамсаргүй дээж сонгож
    загварын таамаглалыг жинхэнэ шошготой харьцуулан үзүүлнэ.
  \item \textbf{Харьцуулах таб:} Flat ба Two-stage архитектурыг нэг дэлгэцэд
    зэрэг ажиллуулж үр дүнг хажуу тийш харьцуулна. Каскад алдааны нөлөөллийг
    (бүлэг~\ref{Chapter6}-д тайлбарласан) шууд ажиглах боломжийг олгоно.
  \item \textbf{Үнэлгээний таб:} Туршилтын олонлог дээрх confusion matrix болон
    ангилал тус бүрийн Precision / Recall / F1 хүснэгтийг статик байдлаар
    харуулна. Загварын нарийвчлалын тоон нотолгоог шууд демоноос харах боломжтой.
\end{enumerate}

Бүх загварыг эхний ачааллах үед санах ойд нэг удаа ачаалж (lazy loading), дараагийн
хүсэлтүүдэд дахин ачаалалгүй хурдан хариу өгдөг. GPU орчинд (AMD DirectML)
инференс нь CPU орчноос мэдэгдэхүйц хурдан бөгөөд DirectML алдаа гарвал систем
автоматаар CPU горимд шилждэг.

\begin{figure}[H]
\centering
\scalebox{0.76}{%
\begin{tikzpicture}[
  wbox/.style={draw, rounded corners=5pt, minimum width=10.0cm,
               minimum height=1.50cm, text width=9.5cm,
               align=center, font=\small, line width=0.8pt},
  sbox/.style={draw, rounded corners=5pt, minimum width=10.0cm,
               minimum height=1.80cm, text width=9.5cm,
               align=center, font=\small, line width=0.8pt,
               fill=purple!8, draw=purple!45},
  bbox/.style={draw, rounded corners=5pt, minimum width=4.6cm,
               minimum height=1.50cm, text width=4.2cm,
               align=center, font=\small, line width=0.8pt},
  decbox/.style={diamond, draw, aspect=2.6, minimum width=3.6cm,
                 inner sep=2pt, align=center, font=\small,
                 fill=yellow!20, draw=yellow!70!black, line width=0.8pt},
  arr/.style={->, >=Stealth, thick, line width=0.9pt},
]

% ── SHARED TOP ───────────────────────────────────────────────────────────
\node[wbox, fill=blue!14, draw=blue!50] (inp) at (0, 0)
  {\textbf{Текст оруулах}\\[3pt]
   {\scriptsize Монгол кирилл · Gradio форм · Нэг эсвэл batch текст}};

\node[wbox, fill=orange!14, draw=orange!55] (pre) at (0,-3.00)
  {\textbf{Урьдчилсан боловсруулалт}\\[3pt]
   {\scriptsize Кирилл шүүлт · нормалчлал · SentencePiece токенчлол}};

\node[sbox] (sel) at (0,-5.60)
  {\textbf{Архитектур сонгох (радио товч)}\\[3pt]
   {\scriptsize \textit{Flat}: нэг шатлалт ($83.26\%$)
    \enspace|\enspace
    \textit{Two-stage}: Stage~1 $\to$ Stage~2 ($78.61\%$)}};

% ── BRANCHES ─────────────────────────────────────────────────────────────
% sel.south = -5.60 - 0.90 = -6.50 ; fork 0.60 below = -7.10
% flat/s1 north = -9.15  (center -9.90, half 0.75)
\node[bbox, fill=green!14, draw=green!55] (flat) at (-2.8,-9.90)
  {\textbf{MN-BERT Flat (4-ангилал)}\\[3pt]
   {\scriptsize max\_len=256 · Эерэг · Саармаг · Хортой сөрөг · Бүтээлч шүүмжлэл}};

\node[bbox, fill=blue!14, draw=blue!50] (s1) at (2.8,-9.90)
  {\textbf{Stage~1: MN-BERT}\\[3pt]
   {\scriptsize 3-ангилал · Эерэг · Саармаг · Сөрөг}};

% Fork from selector
\coordinate (fork) at (0,-7.10);
\draw (sel.south) -- (fork);
\draw[arr] (fork)
           -- node[above, font=\scriptsize, text=green!55!black] {\textit{Flat}}
           (-2.8,-7.10) -- (flat.north);
\draw[arr] (fork)
           -- node[above, font=\scriptsize, text=blue!55!black] {\textit{Two-stage}}
           (2.8,-7.10) -- (s1.north);

% dec center -12.60 ; dec.north = -11.908 ; s1.south = -10.65 ; gap = 1.26 cm
\node[decbox] (dec) at (2.8,-12.60) {Сөрөг уу?};

% s2 center -15.30 ; s2.north = -14.55 ; dec.south = -13.292 ; gap = 1.26 cm
\node[bbox, fill=teal!12, draw=teal!50] (s2) at (2.8,-15.30)
  {\textbf{Stage~2: MN-BERT}\\[3pt]
   {\scriptsize 2-ангилал · Хортой сөрөг · Бүтээлч шүүмжлэл}};

% ── SHARED BOTTOM ────────────────────────────────────────────────────────
% out center -18.10 ; out.north = -17.35 ; s2.south = -16.05 ; gap = 1.30 cm
\node[wbox, fill=red!10, draw=red!45] (out) at (0,-18.10)
  {\textbf{Үр дүн ба итгэлцэл}\\[3pt]
   {\scriptsize Ангиллын шошго · Confidence \%}};

% viz center -20.80 ; viz.north = -20.05 ; out.south = -18.85 ; gap = 1.20 cm
\node[wbox, fill=violet!10, draw=violet!45] (viz) at (0,-20.80)
  {\textbf{Дүрслэл ба тайлан}\\[3pt]
   {\scriptsize Confidence график · Batch CSV татаж авах}};

% ── ARROWS ───────────────────────────────────────────────────────────────
\draw[arr] (inp) -- (pre);
\draw[arr] (pre) -- (sel);

% Two-stage internal
\draw[arr] (s1) -- (dec);
\draw[arr] (dec.south) -- node[right, font=\scriptsize, xshift=3pt] {Тийм} (s2.north);

% "Үгүй": dec.west = (1.0, -12.60) → (0, -12.60) → out.north
\draw[arr] (dec.west) -- node[above, font=\scriptsize] {Үгүй}
           (0,-12.60) -- (out.north);

% Flat → out: descend, bend right, enter box top from above at x=-1.5
\draw[arr] (flat.south) -- (-2.8,-16.50) -- (-1.5,-16.50) -- (-1.5,-17.35);

% Stage 2 → out: descend, bend left, enter box top from above at x=+1.5
\draw[arr] (s2.south) -- (2.8,-16.50) -- (1.5,-16.50) -- (1.5,-17.35);

\draw[arr] (out) -- (viz);

\end{tikzpicture}%
}
\caption{Gradio инференсийн урсгал: нийтлэг урьдчилсан боловсруулалтын дараа
  хэрэглэгч \textit{Flat} (нэг шатлалт, $83.26\%$, max\_len=256) эсвэл
  \textit{Two-stage} (Stage~1 $\to$ нөхцөлт Stage~2, $78.61\%$, max\_len=128)
  архитектурыг сонгоно. Гаралт нь Үр дүн~/ Дүрслэл хэсэгт нэгддэг.}
\label{fig:inference-pipeline}
\end{figure}

%-------------------------------------------------------------------------------
\section{Reproducibility ба туршилтын бүртгэл}

Туршилтын давтагдах байдлыг (reproducibility) хангахын тулд дараах арга хэмжээг
авсан:

\begin{itemize}
  \item \textbf{Random seed тогтворжуулалт:} \texttt{torch.manual\_seed(42)},
    \texttt{numpy.random.seed(42)}, \texttt{random.seed(42)} бүх туршилтанд
    адилхан тохируулсан.
  \item \textbf{Хувилбарын хяналт:} Бүх кодыг Git/GitHub repository-д хадгалж,
    commit hash бүрээр туршилтыг хамтад нь тэмдэглэв. \texttt{requirements.txt}
    файлд бүх сангийн яг тохирох хувилбарыг бичсэн.
  \item \textbf{Туршилтын бүртгэл:} Epoch тутмын loss, accuracy, F1 оноо утгуудыг
    JSON лог файлд хадгалав. Сургалтын муруй (learning curve) болон
    confusion matrix-ыг \texttt{matplotlib}/\texttt{seaborn}-оор хадгалж,
    бүлэг~\ref{Chapter6}-д танилцуулна.
  \item \textbf{Загварын чекпойнт:} Хамгийн сайн validation F1-тэй epoch-ийн
    жинг \texttt{best\_stage1.pt}, \texttt{best\_stage2.pt} файлуудад хадгалав.
  \item \textbf{Орчны тохиргоо:} AMD Radeon RX~9070 GRE 12\,GB GPU дээр
    torch-directml (DirectML backend) ашиглан сургалт явуулсан; нэмэлт
    баталгаажуулалтыг CPU орчинд давтан хийж, үр дүнгийн тогтвортой байдлыг
    шалгав.
\end{itemize}

Эдгээр арга хэмжээ нь guideline-ийн Ш-6 (``reproducible experiment setup
documented'') шаардлагыг бүрэн хангаж байна.

%-------------------------------------------------------------------------------
\section{Бүлгийн дүгнэлт}

Энэ бүлэгт MN-BERT суурьтай Stage~1 ба Stage~2 загваруудын архитектур, сургалтын
дамжлагын бүх алхмууд, гиперпараметрийн сонголт, ангиллын тэнцвэржүүлэлтийн
хоёр арга (SMOTE суурь загварт; Focal Loss + жин $\alpha$ + WeightedRandomSampler
эцсийн нэг шатлалт MN-BERT-д), Gradio-д
суурилсан инференсийн урсгал (Flat $83.26\%$, Two-stage $78.61\%$
гэсэн хоёр архитектур), болон reproducibility-ийн нөхцөлийг
тодорхой баримтжуулсан. Сургалтын бүх бүрэлдэхүүн хэсэг нь нийтлэг Python/PyTorch
экосистем дээр суурилсан тул бусад судлаачид давтан гүйцэтгэх боломжтой. Ийнхүү
загварын хөгжүүлэлтийн үе шат бэлэн болсон бөгөөд дараагийн бүлэгт уг
загваруудын туршилтын үр дүнг нарийвчлан авч үзнэ.
